\documentclass[12pt,a4paper]{report}
\usepackage[T1]{fontenc}
\usepackage[utf8]{inputenc}
\usepackage[french]{babel}
\usepackage[margin=2.5cm]{geometry}
\usepackage{amsmath,amssymb,amsthm}
\usepackage{booktabs}
\usepackage{enumitem}
\usepackage{tikz}
\usepackage{hyperref}
\hypersetup{colorlinks=true,linkcolor=black,citecolor=black,urlcolor=blue,
            pdftitle={Méthodes hybrides pour l'ensemble efficace}}

\newcommand{\Dset}{D}
\newcommand{\E}{E(P_D)}
\newcommand{\Z}{Z}
\newcommand{\fig}[1]{\input{figs/#1}}

\theoremstyle{plain}
\newtheorem{theoreme}{Théorème}[chapter]
\newtheorem{proposition}[theoreme]{Proposition}
\newtheorem{corollaire}[theoreme]{Corollaire}
\theoremstyle{definition}
\newtheorem{definition}[theoreme]{Définition}
\newtheorem{remarque}[theoreme]{Remarque}

\title{\bfseries Méthodes hybrides pour la génération de l'ensemble
des solutions efficaces d'un programme linéaire en nombres entiers
multi-objectif,\\[4pt] avec un critère de décision linéaire fractionnaire}
\author{}
\date{}

\begin{document}
\maketitle

\begin{abstract}
\noindent
Soit $\Dset = \{x \in \mathbb{Z}^n_+ : Ax \leqslant b\}$ non vide et borné, et
soient $\Z_1,\dots,\Z_p$ des critères à maximiser, chacun un rapport de formes
affines de dénominateur strictement positif sur $\Dset$.  Ce mémoire traite
deux questions portant sur l'ensemble efficace $\E$~: le \emph{produire en
entier}, et y maximiser un critère de décision linéaire fractionnaire $\Phi$.

Les méthodes existantes pour la seconde question travaillent dans l'espace des
\emph{décisions}, en tronquant $\Dset$ par des coupes de Sylva--Crema.  Chaque
coupe étant une disjonction, elle coûte $p$ variables binaires et $p+1$
contraintes de type grand-$M$~: chaque sous-problème est strictement plus gros
que le précédent.  Nous conservons la partie inexplorée dans l'espace des
\emph{critères}, sous forme d'une liste de boîtes~: la troncature devient une
soustraction de boîtes, chaque sous-problème est $\Dset$ augmenté de quelques
contraintes linéaires ordinaires, et \textbf{le modèle ne grossit jamais}.

Quatre contributions en découlent.  D'abord une caractérisation qui rend
l'ensemble efficace complet accessible~: \textbf{si $v$ est non dominé, alors
tout $x \in \Dset$ tel que $\Z(x) = v$ est efficace}, de sorte que $\E$ est
exactement la réunion des tranches du front et qu'\emph{aucun test d'efficacité
n'est payé sur un point de tranche}.  Ensuite, une tranche est $\Dset$ augmenté
de $p$ équations \emph{linéaires}, y compris pour des critères fractionnaires~;
résoudre ces équations élimine $r$ variables avant toute énumération et vaut un
facteur médian de $49$.  Ensuite une hybridation qui paie~: une recherche
locale de Pareto fournit des points efficaces vérifiés, placés dans la liste de
boîtes sans sondage, atteignant la totalité de $\E$ à $1{,}36$ fois la méthode
exacte pure sur $120$ instances.  Enfin, un $\E$ explicite admet
\emph{n'importe quelle} fonction de décision calculable, là où les méthodes de
l'espace des décisions exigent une forme de rapport pour construire leurs
coupes.

Sur le problème commun, la recherche dans l'espace des critères l'emporte à
chaque valeur de $p$, de $1{,}52$ à $p=2$ jusqu'à $14{,}57$ à $p=6$, avec des
optima identiques des deux côtés.
\end{abstract}

\tableofcontents

\chapter{Introduction}

\section{Contexte}

En optimisation multi-objectif en nombres entiers, il n'existe en général pas
de solution unique, mais un \emph{ensemble} de solutions efficaces, dont aucune
n'est meilleure que les autres sur tous les critères à la fois.  Le décideur
doit alors trancher, et ce choix s'exprime souvent par un critère
supplémentaire $\Phi$ à optimiser \emph{sur} cet ensemble.

Le problème qui en résulte,
\[
  (P_E) \qquad \max\ \Phi(x) \quad \text{sous} \quad x \in \E ,
\]
est difficile pour une raison structurelle~: $\E$ est une réunion de faces de
$\Dset$, non convexe, et ne possède aucune description explicite par
inégalités.  On ne peut donc pas le passer à un solveur.

\section{Problème traité}

\begin{equation}\label{eq:pd}
  (P_{\Dset}) \qquad \text{<<}\max\text{>>}\ \ \Z_k(x) =
  \frac{c_k^{\top}x + a_k}{d_k^{\top}x + b_k},
  \quad k = 1,\dots,p, \qquad x \in \Dset ,
\end{equation}
avec $\Dset = \{x \in \mathbb{Z}^{n}_{+} : Ax \leqslant b\}$ non vide et borné,
et chaque dénominateur $D_k(x) = d_k^{\top}x + b_k$ strictement positif sur
$\Dset$.  Les deux hypothèses sont \emph{vérifiées} par la mise en {\oe}uvre,
pour une raison développée au chapitre~\ref{ch:valid}~: les violer ne produit
pas une erreur, mais une réponse fausse et d'apparence complète.

\begin{definition}
Un point $x \in \Dset$ est \emph{efficace} lorsqu'il n'existe aucun $y \in
\Dset$ vérifiant $\Z(y) \geqslant \Z(x)$ avec au moins une inégalité stricte.
On note $\E$ l'ensemble des points efficaces et $F = \Z(\E)$ l'ensemble non
dominé.
\end{definition}

\begin{remarque}\label{rem:notZ}
$\Phi$ est une fonction de $x$ et \textbf{non} de $\Z(x)$.  Deux points
efficaces de même vecteur critère peuvent avoir des valeurs de $\Phi$
différentes.  Toute construction menée dans l'espace des critères doit en tenir
compte, et c'est la raison pour laquelle le transfert des méthodes de l'espace
des critères du cas linéaire au cas fractionnaire n'est pas immédiat.
\end{remarque}

\section{Objectifs et plan}

Le travail poursuit les objectifs fixés dans le projet de mémoire~:

\begin{enumerate}[itemsep=3pt]
\item concevoir des méthodes \emph{hybrides} --- exacte--métaheuristique et
      exacte--exacte --- pour l'optimisation d'une fonction linéaire
      fractionnaire sur l'ensemble efficace entier~;
\item produire \emph{un ensemble de solutions efficaces} engendré par ces
      algorithmes~;
\item \textbf{évaluer les performances} des algorithmes proposés par rapport
      aux méthodes existantes.
\end{enumerate}

\noindent
Le chapitre~\ref{ch:etat} situe le travail dans la littérature.  Le
chapitre~\ref{ch:methode} présente la méthode dans l'espace des critères.  Le
chapitre~\ref{ch:contrib} expose les contributions.  Le
chapitre~\ref{ch:resultats} donne l'évaluation des performances, et le
chapitre~\ref{ch:valid} la validation.


\chapter{État de l'art}
\label{ch:etat}

\section{Deux familles de méthodes}

Les méthodes exactes pour l'optimisation sur l'ensemble efficace se répartissent
en deux familles, selon l'espace dans lequel elles cherchent.

\paragraph{L'espace des décisions.}
La région non explorée est maintenue dans l'espace des variables, et réduite
par des coupes.  La coupe de Sylva--Crema supprime $\{x : \Z(x) \leqslant
\Z(\tilde x)\}$ ; comme il s'agit d'une disjonction, elle s'exprime par $p$
variables binaires et $p+1$ contraintes de type grand-$M$.  Toute la
littérature sur un critère \emph{fractionnaire} au-dessus d'un ensemble
efficace entier appartient à cette famille~: Drici, Ouail et Moulaï
(\emph{Annals of Operations Research}, 2018), une contribution de 2015 dans
\emph{RAIRO-Operations Research}, Ouail et Chergui pour le cas hyperbolique,
Chaiblaine, Moulaï et Cherfaoui pour deux fonctions fractionnaires, un article
de 2021 dans \emph{Optimization} pour un objectif quadratique, ainsi que les
variantes stochastiques et à coefficients intervalles de Younsi-Abbaci et
Moulaï.

\paragraph{L'espace des critères.}
La région non explorée est maintenue dans l'espace des valeurs des critères,
sous forme de régions géométriques.  Boland, Charkhgard et Savelsbergh
(\emph{European Journal of Operational Research}, 2017) travaillent dans
l'espace des critères de dimension pleine et le décomposent à chaque itération.
Leurs méthodes \emph{balanced box}, \emph{boxed line} et \emph{triangle
splitting} en sont les précurseurs bi-objectif.  Les synthèses rapportent cette
famille comme actuellement la plus efficace --- \textbf{mais pour un critère de
décision \emph{linéaire}}.

\section{Positionnement de ce travail}

Deux éléments de la construction présentée ici relèvent de l'état de l'art et
ne sont pas revendiqués~: la \textbf{décomposition de l'espace des critères},
établie par les travaux ci-dessus pour un objectif linéaire, et le
\textbf{schéma d'hybridation} métaheuristique--exacte, dont la taxonomie est
connue et dont la variante employée ici --- produire une approximation du front
par une métaheuristique, puis la faire raffiner par une méthode exacte --- est
la forme standard.

Ce que la littérature fractionnaire ne semble pas faire, c'est travailler dans
l'espace des critères.  La combinaison des deux, avec les éléments techniques
qui la rendent \emph{exacte} plutôt qu'approchée, constitue l'apport de ce
mémoire~:

\begin{itemize}[itemsep=3pt]
\item la forme $e_k$ à valeurs entières (section~\ref{sec:ek}), qui remplace
      l'estimation d'un pas minimal par un seuil exact~;
\item les résultats de tranche du chapitre~\ref{ch:contrib}, qui donnent
      l'ensemble efficace complet~;
\item la classe de fonctions de décision admissibles que cet ensemble explicite
      ouvre (section~\ref{sec:admissible}).
\end{itemize}

\begin{remarque}
Cette section indique ce qui a été trouvé et ce qui ne l'a pas été~; elle ne
constitue pas une revendication de nouveauté.  Une revue bibliographique
complète, incluant la lecture intégrale des sources, reste à mener.
\end{remarque}

\chapter{La méthode dans l'espace des critères}
\label{ch:methode}

\section{D'où vient le coût de la méthode en espace des décisions}
\label{sec:cout}

Chaque coupe de Sylva--Crema ajoute $p$ variables binaires et $p+1$ contraintes
au modèle.  Le coût est donc \emph{structurel} et non accidentel~: sur une
instance mesurée, le modèle passe de $5$ colonnes et $7$ lignes à $95$ et
$127$ au fil de $30$ coupes, et le profil montre que $100\%$ du temps est passé
dans le programme en nombres entiers --- la construction des coupes et la
réparation ne comptent pour rien.  Rien en dehors des sous-problèmes ne mérite
d'être optimisé, et c'est précisément leur croissance que la méthode en espace
des critères n'a pas.

\section{La forme \texorpdfstring{$e_k$}{e\_k}, à valeurs entières}
\label{sec:ek}

Les données de chaque critère sont mises à l'échelle entière (numérateur et
dénominateur multipliés par le même facteur, le rapport est inchangé).  On
pose alors
\begin{equation}\label{eq:ek}
  e_k(x ; \bar x) \;=\; D_k(\bar x)\,D_k(x)\,
  \bigl(\Z_k(x) - \Z_k(\bar x)\bigr) ,
\end{equation}
qui est \textbf{linéaire en $x$ et à valeurs entières}, de sorte que
\begin{equation}\label{eq:exact}
  \Z_k(x) > \Z_k(\bar x) \iff e_k(x;\bar x) \geqslant 1 .
\end{equation}

\noindent
C'est le point qui rend le passage des critères linéaires aux critères
fractionnaires \emph{fidèle} et non approché~: aucun pas minimal n'a à être
estimé.  Les deux dénominateurs étant strictement positifs sur $\Dset$,
multiplier par eux ne change ni le sens ni la validité de l'inégalité.

\section{Les boîtes, et pourquoi le modèle ne grossit pas}
\label{sec:boites}

La partie inexplorée est conservée comme une liste de boîtes dans l'espace des
critères.  Tronquer revient à soustraire une boîte, et chaque sous-problème est
$\Dset$ augmenté de quelques lignes linéaires ordinaires issues de
\eqref{eq:exact}.  Aucune variable binaire, aucun grand-$M$, aucune croissance
d'une itération à la suivante.

La partition est disjointe~: l'enfant $k$ reçoit $e_k \geqslant 1$ et
$e_j \leqslant 0$ pour $j < k$, ce qui couvre la boîte privée du centre sans
recouvrement.

\begin{center}
\fig{h-process}
\end{center}


\chapter{Contributions}
\label{ch:contrib}

\section{L'ensemble efficace complet}
\label{sec:complet}

L'énumération du front prouve la complétude de sa liste de \emph{vecteurs},
mais un seul point est conservé par tranche~: la liste de \emph{points} n'est
pas l'ensemble efficace.  Cet écart se comble pour presque rien.

\begin{theoreme}[Caractérisation par tranches]\label{th:tranche}
Si $v$ est non dominé, alors \textbf{tout} $x \in \Dset$ tel que $\Z(x) = v$
est efficace.
\end{theoreme}

\begin{proof}
Supposons qu'un $y \in \Dset$ domine un tel $x$.  Alors $\Z(y) \geqslant \Z(x)
= v$ avec au moins une inégalité stricte, ce qui signifie que $v$ est dominé,
contrairement à l'hypothèse.  Donc aucun $y$ ne domine $x$, et $x$ est efficace.
\end{proof}

\begin{corollaire}
$\E = \{x \in \Dset : \Z(x) = v \text{ pour un } v \text{ non dominé}\}$, et
\textbf{aucun test d'efficacité n'est payé sur un point de tranche}.
\end{corollaire}

\noindent
L'opération la plus coûteuse de la méthode n'est tout simplement pas invoquée.
La démonstration n'utilise aucune propriété de $\Z$ --- ni linéarité, ni
continuité, ni forme de rapport --- de sorte que la caractérisation vaut pour
des critères arbitraires.  Ce qui exige une structure, c'est seulement
l'\emph{énumération} de chaque tranche.

\section{Une tranche est \texorpdfstring{$\Dset$}{D} plus \texorpdfstring{$p$}{p} équations linéaires}
\label{sec:elim}

Fixer $\Z_k(x) = v_k$ est linéaire en $x$ \emph{même pour un critère
fractionnaire}~: comme $d_k^{\top}x + b_k > 0$ sur $\Dset$, l'équation
$(c_k^{\top}x + a_k)/(d_k^{\top}x + b_k) = v_k$ se libère en
\begin{equation}\label{eq:tranche}
  (c_k - v_k d_k)^{\top} x \;=\; v_k b_k - a_k .
\end{equation}

Aucun binaire, aucun grand-$M$, rien qui croisse d'une tranche à la suivante~:
la même propriété que celle de la recherche par boîtes, obtenue une seconde
fois.

\paragraph{Résoudre ces équations plutôt que s'en servir comme coupes.}
L'élimination de Gauss en arithmétique rationnelle exacte exprime $r$ variables
en fonction des autres~; la substitution laisse un système équivalent à $n - r$
inconnues, et le balayage porte sur celui-là.  L'espace de recherche passe de
$(ub+1)^n$ à $(ub+1)^{n-r}$.

\begin{table}[h]
\centering
\caption{Élagage sur les équations de tranche contre résolution de ces
équations.  Quatorze instances, $n = 7 \dots 12$~; ensembles de points
identiques dans tous les cas.}
\begin{tabular}{lrrr}
\toprule
instance & élagage & élimination & rapport \\
\midrule
$n=9$, $p=3$  & $29{,}71$\,s  & $0{,}31$\,s & $96{,}7\times$ \\
$n=10$, $p=5$ & $119{,}47$\,s & $0{,}92$\,s & $129{,}9\times$ \\
$n=12$, $p=3$ & $796{,}84$\,s & $8{,}34$\,s & $95{,}5\times$ \\
\midrule
\textbf{médiane des quatorze} & & & $\mathbf{49{,}2\times}$ \\
\bottomrule
\end{tabular}
\end{table}

\noindent
L'exactitude n'est pas affectée~: tous les coefficients sont des rationnels
exacts et l'élimination l'est aussi, de sorte qu'il s'agit d'un changement
d'ordre de parcours et non d'une approximation.

\section{Ce qu'un \texorpdfstring{$\E$}{E(P\_D)} explicite autorise}
\label{sec:admissible}

C'est la contribution que l'on défendrait en premier, parce qu'elle ne repose
pas sur le fait d'être le plus rapide~: elle dit quels problèmes peuvent être
\emph{posés}.

\begin{quote}
Une méthode en espace des décisions exige une \emph{structure dans $\Phi$} et
rien de $\E$~; cette voie exige \emph{$\E$} et rien de $\Phi$.
\end{quote}

\begin{theoreme}[Classe admissible]\label{th:admissible}
Soit $\E$ connu et prouvé complet.  Alors pour \emph{toute} fonction $f$
définie sur $\Dset$, $\max\{f(x) : x \in \E\}$ et l'ensemble complet de ses
maximiseurs se calculent exactement en évaluant $f$ sur $\E$.  Aucune
linéarité, monotonie, continuité, dérivabilité ni forme de rapport n'est
requise.
\end{theoreme}

\noindent
La démonstration est la finitude de $\E$ et rien d'autre.  Le contenu n'est pas
dans l'énoncé~: il est dans le fait d'avoir payé $\E$ à un prix qui rend ce
paiement intéressant, et l'énoncé est ce que ce prix achète.

\begin{table}[h]
\centering
\caption{Six formes de fonction de décision sur l'instance publiée, chacune
vérifiée contre une maximisation exhaustive de référence.  $\Phi_1,\Phi_2$ sont
linéaires fractionnaires.}
\begin{tabular}{llr}
\toprule
fonction de décision & admise par un branch-and-cut ? & optimum \\
\midrule
$\Phi_1$                        & oui --- le cas classique & $5/17$ \\
$\Phi_1 + \Phi_2$               & non & $47/35$ \\
$\max(\Phi_1,\Phi_2)$           & non & $6/5$ \\
$\Phi_1 \cdot \Phi_2$           & non & $6/35$ \\
$(x_1^2+1)/(x_2^2+1)$           & non & $10$ \\
définie par morceaux            & non & $3/11$ \\
\bottomrule
\end{tabular}
\end{table}

\noindent
Un branch-and-cut pour un $\Phi$ linéaire fractionnaire tire ses coupes et ses
bornes de la forme de rapport.  Remplacer $\Phi$ par une \emph{somme} de deux
rapports ne le ralentit pas~: il cesse de s'appliquer.  Ici, ce cas coûte
exactement ce que coûte un rapport unique.

\paragraph{Plusieurs fonctions de décision à la fois.}
Optimiser \emph{deux} fonctions linéaires fractionnaires sur un ensemble
efficace entier constitue un problème à part entière dans la littérature, avec
son propre branch-and-bound.  Ici, c'est un filtre de dominance sur un ensemble
déjà en main, et le coût ne dépend ni du nombre de fonctions ni de leur nature.
Sur l'instance publiée, le filtre à deux rapports conserve $4$ des $7$ points
efficaces, exactement la réponse exhaustive.

\section{L'hybridation qui paie}
\label{sec:hybride}

Une recherche locale de Pareto avec archive fournit des points efficaces
\textbf{vérifiés}, placés dans la liste de boîtes \textbf{sans sondage}.  La
correction se situe là~: l'archive est filtrée par dominance \emph{parmi ce
qu'elle a vu}, ce qui n'est pas le test exact~; rien ne franchit la frontière
sans vérification.  La métaheuristique décide \emph{où chercher}~; elle ne
décide jamais de ce qui est vrai.

\paragraph{L'ordre est toute la différence.}
C'est le résultat pratique le plus transférable du travail.  L'hybridation ne
paie que si les graines sont transmises dans l'ordre où le sondage de la
méthode exacte les aurait rencontrées.

\begin{table}[h]
\centering
\caption{Sondages contre l'énumération non amorcée, selon l'ordre de
transmission.  Vingt-sept instances, aucune censurée~; le nombre de sondages
est déterministe.}
\begin{tabular}{lrll}
\toprule
ordre & sondages & IQR & a augmenté le compte sur \\
\midrule
ordre de l'archive & $0{,}83\times$ & $[0{,}74,\ 0{,}92]$ & $2$ sur $27$ \\
\textbf{ordre du sondage} & $\mathbf{0{,}69}\times$ & $[0{,}67,\ 0{,}75]$ & \textbf{aucune} \\
inversé (témoin) & $0{,}90\times$ & $[0{,}83,\ 1{,}05]$ & $9$ sur $27$ \\
\bottomrule
\end{tabular}
\end{table}

\noindent
Le bras inversé est le témoin~: il est \emph{tout aussi trié} et ne diffère de
l'ordre du sondage que par la direction.  Or l'ordre du sondage n'augmente le
nombre de sondages sur aucune des vingt-sept instances, tandis que l'inversé
l'augmente sur neuf, et sa médiane se place avec l'ordre de l'archive.  Ce que
l'amorçage exploite est donc la \emph{direction propre de la recherche}, et non
le fait qu'une liste soit triée.

\paragraph{Et plus précisément~: la première graine.}
La pré-partition consomme les graines \emph{dans l'ordre}, de sorte que la
première fixe la structure de boîtes de premier niveau.  En conditionnant sur
elle, sur $228$ ordres et $26$ instances~: lorsque la première graine est celle
que le sondage aurait rencontrée en premier, la médiane est $0{,}708$~; sinon
$0{,}833$.  Cet écart représente l'essentiel de l'écart entre les deux ordres
nommés.


\chapter{Évaluation des performances}
\label{ch:resultats}

Ce chapitre répond au quatrième point du projet~: l'évaluation des performances
des algorithmes proposés par rapport aux méthodes existantes.

\section{Protocole de mesure}

Toute l'arithmétique est rationnelle exacte, sans aucun flottant, de sorte que
chaque comparaison oppose deux problèmes résolus exactement.  La statistique
est la \textbf{médiane des rapports par instance} avec l'écart interquartile,
jamais un rapport de totaux --- lequel serait décidé par la seule instance la
plus lourde.

\paragraph{Ce qu'est la référence.}
Le côté « espace des décisions » de chaque comparaison est une
\textbf{réimplémentation} de la méthode publiée, sur des instances engendrées
ici --- et non le code de ses auteurs ni leurs instances de référence.  Les
deux côtés partagent le même simplexe, la même arithmétique rationnelle exacte,
le même branch-and-bound et la même machine, et la réimplémentation suit
l'algorithme imprimé pas à pas.  La comparaison \emph{isole donc la différence
algorithmique}, ce qui est une question plus propre qu'une comparaison de code
à code --- celle-ci mesurerait aussi le langage, le solveur sous-jacent et
l'effort d'ingénierie.

En contrepartie, \textbf{les secondes absolues ne portent aucune information
hors de cette mise en {\oe}uvre}~: du Python pur sur des rationnels exacts est
bien plus lent que tout solveur de production, également des deux côtés.  Seuls
les rapports se transportent, et seuls des rapports sont cités.

\section{Espace des décisions contre espace des critères}

\begin{table}[h]
\centering
\caption{Le gain selon le nombre de critères, $n = 6$, \textbf{trente
instances par ligne}.}
\label{tab:marge}
\begin{tabular}{crrrr}
\toprule
$p$ & décisions & critères & rapport (médiane) & IQR \\
\midrule
2 & $0{,}07$\,s & $0{,}04$\,s & $1{,}52\times$  & $[1{,}25,\ 1{,}91]$ \\
3 & $0{,}34$\,s & $0{,}12$\,s & $2{,}69\times$  & $[1{,}87,\ 3{,}72]$ \\
4 & $0{,}50$\,s & $0{,}13$\,s & $4{,}99\times$  & $[2{,}39,\ 11{,}62]$ \\
5 & $2{,}19$\,s & $0{,}25$\,s & $7{,}93\times$  & $[3{,}61,\ 19{,}98]$ \\
6 & $3{,}67$\,s & $0{,}31$\,s & $\mathbf{14{,}57\times}$ & $[7{,}90,\ 25{,}17]$ \\
\bottomrule
\end{tabular}
\end{table}

\begin{center}
\fig{r-campaign-margin}
\end{center}

\noindent
\textbf{La méthode l'emporte à chaque ligne, et l'avantage croît avec $p$} ---
la direction qui importe, puisque c'est $p$ qui rend ces problèmes difficiles.
Trois des $150$ exécutions ont épuisé un budget de $45$\,s et sont exclues
plutôt que comptées à leur budget~; cette censure est \emph{unilatérale} --- 
elle ne peut que faire paraître la méthode en espace des décisions meilleure ---
de sorte que chaque marge citée est conservatrice.  Toute exécution terminée a
donné le même optimum des deux côtés.

\paragraph{La référence est-elle handicapée ?}
Le perdant a été profilé plutôt que simplement cité.  Sur l'instance qui
demande $137$\,s pour produire un front de $30$ vecteurs en espace des
décisions, $100\%$ du temps est dans le programme en nombres entiers, et le
modèle croît de $5$ colonnes à $95$.  Le coût est structurel.  Le seul
handicap plausible --- sonder en maximisant une direction plutôt qu'avec un
objectif nul --- a été testé~: l'objectif nul est $4$ à $13$ fois
\emph{pire}, à coupes identiques.

\section{L'hybridation}

\begin{table}[h]
\centering
\caption{L'hybride contre la méthode exacte pure. Réponses identiques partout.}
\begin{tabular}{lrr}
\toprule
expérience & instances & rapport \\
\midrule
$\E$ complet, $n = 6\dots12$ & $120$ & $\mathbf{1{,}36\times}$, plus rapide sur $116$ \\
critères véritablement fractionnaires & $30$ & $1{,}39\times$ \\
cinq familles structurelles & $150$ & $1{,}23\times$ médian, plus rapide sur $136$ \\
\bottomrule
\end{tabular}
\end{table}

\noindent
Chacune des cinq familles paie, entre $1{,}14\times$ et $1{,}37\times$~:
l'avantage n'est donc pas un artefact d'une famille particulière.

\section{Le domaine de fonctionnement}

Deux quantités gouvernent le coût, et elles sont distinctes.  L'énumération du
front suit $|F|$~; le complètement par tranches suit $n - p$, parce que les
équations de tranche sont de rang $p$ --- mesuré par instance, et égal à $p$
dans chaque cas --- de sorte que l'élimination laisse $n-p$ variables libres.
Le temps de tranche le confirme~: $0{,}65$, $2{,}09$, $48{,}08$ et $126{,}14$
secondes à $n-p = 7, 9, 11, 13$.

À $p=3$, la méthode est confortable jusqu'à $n=12$ et résout encore des
instances à $n=14$--$16$.  \textbf{Toutes les instances présentées dans ce
mémoire se situent à $n \leqslant 12$}, à l'intérieur de ce domaine avec de la
marge~: chaque chiffre rapporté provient donc du régime où la méthode
fonctionne à l'écart des deux limites.

\section{Bilan des essais}

Vingt-deux modifications ont été essayées et mesurées.  Les plus instructives~:

\begin{table}[h]
\centering
\caption{Essais et verdicts.}
\begin{tabular}{lll}
\toprule
essai & verdict & effet \\
\midrule
espace des décisions $\to$ critères & \textbf{gain} & $1{,}5\times$ à $14{,}6\times$ \\
amorçage par marche de Pareto & \textbf{gain} & $\mathbf{1{,}31\times}$ vers tout $\E$ \\
résolution des équations de tranche & \textbf{gain} & $49\times$ médian \\
tranches au lieu d'un test d'efficacité & \textbf{gain} & tout $\E$, sous $2\times$ \\
dimensionnement de la marche & \textbf{gain} & $1{,}00 \to 1{,}53\times$ à $n=12$ \\
hybride sur critères fractionnaires & \textbf{gain} & $1{,}39\times$ \\
Tchebychev comme générateur & \textbf{gain} & atteint les points non supportés \\
démarrage à chaud + point mémorisé & \textbf{gain} & $20\%$ de pivots en moins \\
\bottomrule
\end{tabular}
\end{table}


\chapter{Validation}
\label{ch:valid}

\section{Comment les résultats sont contrôlés}

La mise en {\oe}uvre est en Python pur, sans aucune dépendance externe, en
arithmétique rationnelle exacte de bout en bout --- aucun flottant nulle part,
donc aucune tolérance nulle part --- et porte \textbf{112 tests} exécutés à
chaque validation.

Plusieurs tests contrôlent des \emph{propriétés} plutôt que des réponses, car
une méthode qui perdrait silencieusement des points efficaces s'accorderait
quand même avec l'optimum sur les instances favorables~: la partition est
disjointe~; la métaheuristique ne transmet jamais de point inefficace~; tout
point mémorisé est efficace~; une racine démarrée à chaud atteint la même
valeur que la racine froide~; la sortie anticipée laisse des boîtes ouvertes et
prouve tout de même l'optimalité.

Les réponses sont recoupées de \textbf{quatre façons indépendantes}~:
l'énumération exhaustive, une énumération de l'ensemble non dominé par coupes
répétées, un balayage trié avec tests de dominance, et un certificat
d'optimalité qui règle $362\,293$ concurrents avec $3$ tests d'efficacité
exacts.

\section{Les hypothèses sont vérifiées, pas supposées}

Les deux hypothèses de \eqref{eq:pd} --- $\Dset$ borné et dénominateurs
strictement positifs --- sont \textbf{vérifiées} par les points d'entrée, pour
une raison précise~: les violer ne produit pas une erreur mais une réponse
fausse d'apparence complète.  Sur un $\Dset$ non borné pourtant admissible,
l'énumération rendait un \emph{front vide marqué complet}~; sur un critère dont
le dénominateur change de signe, elle rendait un seul vecteur comme front
entier alors qu'un point admissible le domine.

Le coût de ces vérifications est une charge fixe par problème~: un programme en
nombres entiers pour la bornitude, un par critère \emph{véritablement}
fractionnaire pour les dénominateurs et aucun pour un critère linéaire, et $n$
programmes mis en cache pour la boîte de variables.  Face aux volumes de
sous-problèmes du chapitre~\ref{ch:resultats}, ils ne se voient pas.

\section{Ce que la démarche a permis de corriger}

La discipline de mesure adoptée a permis de corriger des énoncés en cours de
route, ce qui est rapporté ici parce que cela indique la fiabilité des chiffres
conservés~:

\begin{itemize}[itemsep=3pt]
\item la statistique est la médiane des rapports \emph{par instance} --- un
      rapport de totaux avait un jour été dominé par une seule instance
      fournissant $110{,}95$ des $112{,}77$ secondes~;
\item le \textbf{travail déterministe} est compté à côté des secondes ---
      sous-programmes, n{\oe}uds de branch-and-bound, pivots --- car le plancher
      de bruit de la machine est de $\pm 30\%$ par instance et aucun nombre de
      répétitions ne le supprime~;
\item une corrélation calculée sur des médianes par classe donnait $-0{,}98$
      là où les paires brutes donnent $-0{,}52$~: c'est le chiffre brut qui est
      cité.
\end{itemize}

\chapter{Conclusion et perspectives}

\section{Ce qui a été obtenu}

\begin{enumerate}[itemsep=4pt]
\item \textbf{L'ensemble efficace complet $\E$}, et non le seul front, grâce au
      théorème~\ref{th:tranche} et sans payer de test d'efficacité sur un point
      de tranche.  C'est ce que le projet attendait~: un ensemble de solutions
      efficaces engendré par les algorithmes proposés --- et il est livré
      \emph{complet} là où le projet n'en demandait qu'un sous-ensemble.
\item \textbf{Une hybridation exacte--métaheuristique qui paie}~: $1{,}36\times$
      sur $120$ instances, plus rapide sur $116$, $1{,}39\times$ sur des
      critères fractionnaires, et un gain dans chacune des cinq familles.
\item \textbf{L'élimination des variables de tranche}~: le plus grand gain
      unitaire du travail, $49\times$ médian.
\item \textbf{La classe de décision admissible} (théorème~\ref{th:admissible})~:
      un $\E$ explicite accepte toute fonction calculable, là où les méthodes
      existantes exigent une forme de rapport.  C'est le résultat que l'on
      défendrait en premier, car il ne porte pas sur la vitesse mais sur les
      problèmes que l'on peut poser.
\item \textbf{L'évaluation des performances} demandée par le projet~: de
      $1{,}52\times$ à $14{,}57\times$ selon $p$, avec une censure unilatérale
      en faveur de la méthode comparée et des optima identiques.
\end{enumerate}

\section{Perspectives}

\begin{itemize}[itemsep=3pt]
\item \textbf{Comparaison avec les mises en {\oe}uvre des auteurs} sur leurs
      instances de référence.  La comparaison menée ici isole la différence
      algorithmique~; celle-là mesurerait l'écart de bout en bout.
\item \textbf{Élargissement du domaine de fonctionnement.}  Les deux quantités
      qui le gouvernent étant identifiées --- $|F|$ et $n-p$ --- une réduction
      du coût de l'énumération du front élargirait directement le domaine.
\item \textbf{Critères à ensembles de niveau non affines.}  La machinerie de
      tranche s'étend à toute réunion finie de polyèdres, ce qui couvre déjà
      $|c^{\top}x+a|/(d^{\top}x+b)$ et les critères minimax~; l'énumération du
      front reste à étendre pour que ces classes soient traitées de bout en
      bout.
\item \textbf{Revue bibliographique complète}, incluant la lecture intégrale
      des sources citées au chapitre~\ref{ch:etat}.
\end{itemize}

\end{document}
